\chapter{Die Konfigurationsdatei}\label{ch:konfiguration}

Jede Rechnung wird durch eine Konfiguration im JSON-Format beschrieben. \cmd{fgsolve} liest sie aus einer Datei,
Python nimmt sie als \texttt{dict}, die GUI bearbeitet sie in Formularen. Es ist in allen drei Fällen
\emph{dasselbe Format} mit demselben Parser.

\section{Gesamtaufbau}
\begin{lstlisting}[language=json]
{
  "dimension": "1D",              // Geometrie: "1D" oder "2D"
  "solver": "rcwa",               // "rcwa" oder "rk"
  "truncation": {...},            // Fourierordnungen
  "period_um": {...},             // Perioden
  "incidence": {...},             // Einfall, auch als Scan
  "sweep": {...},                 // optional: Parameter-Sweep
  "materials": {...},             // Materialtabelle
  "superstrate": "air",           // Einfallsmedium (oben)
  "substrate": "glass",           // Substrat (unten)
  "layers": [...],                // Schichten von unten nach oben
  "rk": {...}, "rcwa": {...},     // Löseroptionen
  "nearfield": {...},             // optional: Nahfeld
  "output": {...}                 // Ausgabe
  // weitere optionale Blöcke: Leistung, Genauigkeit, ASR, PML, Zylinder (Kapitel 11, 15)
}
\end{lstlisting}

\begin{hinweis}[{Groß- und Kleinschreibung, Reihenfolge, unbekannte Schlüssel}]
Schlüssel sind klein geschrieben. Die Reihenfolge der Schlüssel spielt keine Rolle, die Reihenfolge der Schichten
dagegen schon. Unbekannte Schlüssel werden ignoriert; ein Tippfehler wie \key{thicknes\_um} führt deshalb meist zu
der Fehlermeldung, dass ein \emph{Pflicht}schlüssel fehlt. Der Schlüssel \key{reference} in den Beispielen enthält
Sollwerte für die Tests und hat auf die Rechnung keinen Einfluss.
\end{hinweis}

\section{Grundeinstellungen}
\begin{center}
\small
\begin{tabularx}{\linewidth}{lllY}
\toprule
Schlüssel & Typ & Standard & Bedeutung \\
\midrule
\key{dimension} & Text & Pflicht & \val{"1D"} (periodisch in $x$) oder \val{"2D"} (periodisch in $x$ und $z$) \\
\key{solver} & Text & \val{"rcwa"} & \val{"rcwa"} oder \val{"rk"} (Kapitel~\ref{ch:loeser}) \\
\key{truncation.Nx} & Zahl & 10 & höchste Fourierordnung in $x$ \\
\key{truncation.Nz} & Zahl & 0 & höchste Fourierordnung in $z$ (nur 2D, sonst 0) \\
\key{period\_um.x} & Zahl & Pflicht & Periode $d_x$ in \textmu m \\
\key{period\_um.z} & Zahl & 0 & Periode $d_z$; bei \val{"2D"} Pflicht und $>0$ \\
\bottomrule
\end{tabularx}
\end{center}
Der Block \key{truncation} selbst ist Pflicht (auch wenn er nur \val{\{"Nx": 15\}} enthält). Bei
rotationssymmetrischen Strukturen (Abschnitt~\ref{sec:zylinder}) entfallen \key{dimension}, \key{truncation} und
\key{period\_um}.

\section{Materialien}\label{sec:materialien}
\index{Material}\index{Anisotropie}\index{Euler-Winkel}%
Der Block \key{materials} ordnet Namen ihre optischen Eigenschaften zu. Alle anderen Stellen der Konfiguration
verweisen nur noch auf diese Namen.
\begin{lstlisting}[language=json]
"materials": {
  "air":  { "n": 1.0 },                         // reelle Brechzahl
  "au":   { "n": [0.2, 3.4] },                  // komplexe Brechzahl n + ik
  "ag":   { "eps": [-15.9, 1.1] },              // Permittivität statt Brechzahl
  "co":   { "eps_tensor": [[[-8.19,16.38],[0,0],[-0.495,-0.106]],
                           [[0,0],[-8.19,16.38],[0,0]],
                           [[0.495,0.106],[0,0],[-8.19,16.38]]] },
  "lnb":  { "eps_principal": [4.98, 4.98, 4.6], "euler_deg": [0, 90, 0] }
}
\end{lstlisting}
\begin{center}
\small
\begin{tabularx}{\linewidth}{lY}
\toprule
Angabe & Bedeutung \\
\midrule
\key{n} & Brechzahl, reell (\val{1.5}) oder komplex als \val{[Re, Im]}; intern $\varepsilon=n^2$ \\
\key{eps} & relative Permittivität, reell oder \val{[Re, Im]} \\
\key{eps\_tensor} & voller $3\times3$-Tensor in der Reihenfolge $(x,y,z)$; jeder Eintrag \val{[Re, Im]}. Zeile $i$,
Spalte $j$ ist $\varepsilon_{ij}$ \\
\key{eps\_principal} & drei Hauptwerte (reell oder \val{[Re, Im]}) und \key{euler\_deg}: Drehung $R$ als
Euler-Winkel in der Folge $z$-$x$-$z$; $\varepsilon = R\,\mathrm{diag}(\varepsilon_1,\varepsilon_2,\varepsilon_3)\,R^T$ \\
\bottomrule
\end{tabularx}
\end{center}
Jedes Material braucht genau eine dieser Angaben. Fehlt sie, meldet \FG{}
\texttt{material 'x' needs 'n', 'eps', 'eps\_tensor' or 'eps\_principal'}.

\begin{achtung}[Vorzeichen]
Absorption heißt $\operatorname{Im}n>0$ (Zeitkonvention $e^{-i\omega t}$). Tabellenwerte aus Quellen mit der
Konvention $e^{+i\omega t}$ (dort $n-ik$) müssen Sie mit positivem Imaginärteil eintragen.
\end{achtung}

\paragraph{Superstrat und Substrat.} \key{superstrate} und \key{substrate} nennen je einen Materialnamen. Das
Superstrat muss reell und isotrop sein (das Licht kommt von dort). Das Substrat darf absorbieren; dann gibt es
keine transmittierten Ordnungen. Anisotrope Substrate sind erlaubt.

\subsection{Dispersive Materialien}\label{sec:conf-dispersion}
\index{Dispersion}\index{Sellmeier}\index{Cauchy}\index{Drude--Lorentz}\index{Materialtabelle}%
Eine feste Zahl gilt für alle Wellenlängen eines Scans. Für Spektren über mehr als einige Nanometer ist das bei
Metallen und Gläsern falsch. Deshalb darf jeder Materialwert unter \key{n}, \key{eps} und in jedem Eintrag von
\key{eps\_principal} statt einer Zahl ein \emph{Modell} sein. \FG{} wertet es an jedem Scanpunkt bei dessen
Wellenlänge aus; die Geometrie bleibt fest.
\begin{lstlisting}[language=json]
"materials": {
  "sio2": { "n": { "sellmeier": { "A": 1.0, "B": [0.6961663, 0.4079426, 0.8974794],
                                  "C_um2": [0.00467914826, 0.0135120631, 97.9340025],
                                  "range_um": [0.21, 6.7] } } },   // optional
  "tio2": { "n": { "cauchy": { "A": 2.40, "B_um2": 0.03, "C_um4": 0.0, "k": 0.0 } } },
  "au":   { "n": { "table": { "file": "materials/au.csv", "unit": "um",
                              "interpolation": "cubic", "extrapolate": "error" } } },
  "ag":   { "eps": { "drude_lorentz": { "eps_inf": 1.0, "omega_p_eV": 9.01,
              "f0": 0.845, "gamma0_eV": 0.048,
              "oscillators": [ { "f": 0.065, "omega_eV": 0.816, "gamma_eV": 3.886 },
                               { "f": 0.124, "omega_eV": 4.481, "gamma_eV": 0.452 },
                               { "f": 0.011, "omega_eV": 8.185, "gamma_eV": 0.065 },
                               { "f": 0.840, "omega_eV": 9.083, "gamma_eV": 0.916 },
                               { "f": 5.646, "omega_eV": 20.29, "gamma_eV": 2.419 } ] } } },
  "lnb":  { "eps_principal": [ { "sellmeier": {...} }, { "sellmeier": {...} }, 4.6 ],
            "euler_deg": [0, 90, 0] }
}
\end{lstlisting}
\begin{center}
\small
\begin{tabularx}{\linewidth}{llY}
\toprule
Modell & liefert & Formel und Schlüssel \\
\midrule
\key{sellmeier} & $\varepsilon$ & $\varepsilon=A+\sum_k B_k\lambda^2/(\lambda^2-C_k)$, $C_k$ in \textmu m$^2$
(\key{C\_um2}), \key{A} Standard 1, reell. Optional \key{range\_um}: außerhalb Fehler \\
\key{cauchy} & $n$ & $n=A+B/\lambda^2+C/\lambda^4+ik$ (\key{B\_um2}, \key{C\_um4}, \key{k}, Standard 0) \\
\key{drude\_lorentz} & $\varepsilon$ & $\varepsilon=\varepsilon_\infty-\dfrac{f_0\omega_p^2}{\omega(\omega+i\Gamma_0)}
+\sum_j\dfrac{f_j\omega_p^2}{\omega_j^2-\omega^2-i\omega\Gamma_j}$, Energien in eV,
$\hbar\omega=1{,}239841984/\lambda$ \\
\key{table} & Schlüssel & unter \key{n}: Spalten $\lambda$, $n$[, $k$]; unter \key{eps}/\key{eps\_principal}:
$\lambda$, $\operatorname{Re}\varepsilon$[, $\operatorname{Im}\varepsilon$] \\
\bottomrule
\end{tabularx}
\end{center}
\begin{itemize}
  \item Liefert ein Modell $n$, gilt $\varepsilon=n^2$; liefert es $\varepsilon$, gilt das auch unter \key{n}.
        Verluste heißen $\operatorname{Im}\varepsilon>0$ bzw.\ $k\ge0$.
  \item \textbf{Tabellen}: \key{file} (CSV; Trennzeichen Komma, Semikolon oder Leerraum; Zeilen mit \texttt{\#} und
        eine Kopfzeile ohne Zahlen werden übersprungen; relativ zum Verzeichnis der Konfiguration) \emph{oder}
        \key{data} mit den Zeilen. \key{unit} der ersten Spalte: \val{um} (Standard), \val{nm}, \val{eV}.
        \key{interpolation}: \val{linear} (Standard) oder \val{cubic} (Not-a-knot-Spline, ab 4 Zeilen); jede
        Spalte für sich in $\lambda$. Außerhalb der Tabelle ist die Auswertung ein Fehler, mit
        \val{"extrapolate": "constant"} gilt der Randwert mit einer Warnung in der Diagnose.
  \item Dispersive Einträge in \key{eps\_tensor} gibt es nicht; dafür \key{eps\_principal} mit \key{euler\_deg}.
  \item \FG{} prüft beim Laden jeden Scanpunkt und jeden Sweep-Wert: Gültigkeitsbereich, Sellmeier-Pol und ein
        verlustfreies Superstrat. Die Meldung nennt den JSON-Pointer, z.\,B.\
        \texttt{/materials/au/n/table: lambda 0.25 um outside [0.3, 1.9] um}.
  \item Der Einfallswinkel $\varphi$ gilt im Superstrat. Ist dessen Brechzahl dispersiv, hängt die
        Tangentialkomponente $\alpha_0=\nu_M(\lambda)\sin\varphi\cos\psi$ von $\lambda$ ab (fester Winkel, nicht
        fester Wellenvektor).
  \item Modellwerte sind Zahlen der Konfiguration: Sweeps, \texttt{set\_parameter} und Gradienten erreichen sie per
        Pointer, z.\,B.\ \val{/materials/sio2/n/\allowbreak sellmeier/B/0} oder \val{/materials/ag/eps/\allowbreak table/data/3/2}
        (Tabellenwerte nur mit \key{data}; die Wellenlängenspalte ist kein Gradientenparameter).
  \item Ein Modell mit konstantem Wert ergibt bitgleich dasselbe wie die Zahl.
\end{itemize}
\begin{hinweis}[Herkunft der Werte]
Die Drude--Lorentz-Werte für Silber oben (und für Aluminium in \file{docs/config\_format.md}) stammen von Raki\'c
et al.\ \ref{lit:rakic} und sind gegen die Tabellen von refractiveindex.info geprüft. Raki\'c hat an Messdaten mehrerer Quellen
angepasst: Gegenüber Johnson--Christy \ref{lit:jc} ist $|\operatorname{Re}\varepsilon|$ für Silber im Sichtbaren etwa 20\,\%
kleiner und $\operatorname{Im}\varepsilon$ zwei- bis fünfmal größer. Tabellen übernehmen Sie aus refractiveindex.info
mit \texttt{fouriergrating.materials.from\_refractiveindex} (Abschnitt~\ref{sec:py-dispersion}).
\end{hinweis}

\section{Einfall}\label{sec:conf-einfall}
\index{Einfall}%
\begin{center}
\small
\begin{tabularx}{\linewidth}{lllY}
\toprule
Schlüssel & Typ & Standard & Bedeutung \\
\midrule
\key{lambda\_um} & Wert & Pflicht & Vakuumwellenlänge in \textmu m \\
\key{phi\_deg} & Wert & 0 & Polarwinkel gegen die Normale \\
\key{psi\_deg} & Wert & 0 & Azimut gegen $x$ \\
\key{delta\_deg} & Wert & 0 & Polarisation: 0 = s (TE), 90 = p (TM) \\
\key{type} & Text & -- & nur für Sonderfälle: \val{"gaussian\_beam"}, \val{"plane\_wave"} (Zylinder) \\
\bottomrule
\end{tabularx}
\end{center}
Jeder \emph{Wert} darf sein:
\begin{itemize}
  \item eine Zahl: \val{0.6328},
  \item eine Liste: \val{[0.55, 0.6328, 0.7]},
  \item ein Bereich: \val{\{"from": 0.5, "to": 0.8, "steps": 31\}}. Die Endpunkte sind eingeschlossen, die Werte
        gleichabständig; \key{steps} $=1$ ergibt nur \key{from}.
\end{itemize}
Gerechnet wird das \textbf{kartesische Produkt} $\lambda\times\varphi\times\psi\times\delta$ in dieser
Schachtelung ($\lambda$ außen, $\delta$ innen). Zwei Wellenlängen und beide Polarisationen
(\val{"delta\_deg": [0, 90]}) ergeben also vier Rechnungen. Einzelheiten in Kapitel~\ref{ch:scans}.

\section{Sweep}
\index{Sweep}%
Ein Sweep variiert eine beliebige \emph{Zahl} der Konfiguration, adressiert durch einen JSON-Pointer:
\begin{lstlisting}[language=json]
"sweep": { "parameter": "/layers/0/thickness_um",
           "values": { "from": 0.05, "to": 1.2, "steps": 24 } }
\end{lstlisting}
\key{values} hat dieselbe Form wie ein Einfallswert. Der Pointer muss auf eine vorhandene Zahl zeigen. Für jeden
Sweep-Wert wird der ganze Scan gerechnet. Kapitel~\ref{ch:scans} erklärt JSON-Pointer ausführlich.

\section{Schichten}
\key{layers} ist eine Liste von Schichten \emph{von unten nach oben}; jede hat einen \key{type}. Eine leere Liste
(oder ein fehlender Block) beschreibt eine einfache Grenzfläche zwischen Substrat und Superstrat. Die Schichttypen
beschreibt Kapitel~\ref{ch:schichten} ausführlich; hier die Übersicht:
\begin{center}
\small
\begin{tabularx}{\linewidth}{lYcc}
\toprule
\key{type} & Beschreibung & 1D & 2D \\
\midrule
\val{homogeneous} & homogene Schicht & \checkmark & \checkmark \\
\val{lamellar1d} & Lamellen, Materialwechsel an Stellen in $x$ & \checkmark & \checkmark ($z$-invariant) \\
\val{profile1d} & Grenzfläche $y=g(x)$: Sinus, Dreieck, Spline & \checkmark & \checkmark ($z$-invariant) \\
\val{polygon2d} & Querschnitt aus Polygonen und Rechtecken & -- & \checkmark \\
\val{contour2d} & Querschnitt aus Kreisen, Ellipsen, Polygonen, NURBS, mit Löchern & -- & \checkmark \\
\val{profile2d} & Grenzfläche $y=g(x,z)$: Eierkarton, Sinus in $x$ & -- & \checkmark \\
\val{cad} & STEP-Körper, beim Laden in Bänder zerlegt & -- & \checkmark \\
\val{rings} & konzentrische Ringe (Zylinderpfad) & \multicolumn{2}{c}{nur Zylinder} \\
\bottomrule
\end{tabularx}
\end{center}

\section{Löseroptionen}
\index{Löseroptionen}%
\begin{center}
\small
\begin{tabularx}{\linewidth}{lllY}
\toprule
Schlüssel & Typ & Standard & Bedeutung \\
\midrule
\key{rk.tableau} & Text & \val{"gauss2"} & Runge--Kutta-Verfahren: \val{gauss2}, \val{gauss3}, \val{radau3} \\
\key{rk.tol} & Zahl & $10^{-8}$ & Toleranz der Schrittweitensteuerung \\
\key{rcwa.slices} & Zahl & 40 & Treppenstufen je Profilschicht (RCWA); je Schicht überschreibbar \\
\key{backend} & Text & \val{"auto"} & \val{cpu}, \val{cuda}, \val{auto} (Abschnitt~\ref{sec:gpu}) \\
\key{backend\_threshold} & Zahl & 160 & Matrixgröße, ab der \val{auto} die GPU nutzt \\
\key{precision} & Text & \val{"auto"} & \val{fp64}, \val{mixed}, \val{fp32}, \val{auto} (Abschnitt~\ref{sec:genauigkeit}) \\
\key{precision\_fallback} & bool & auto & bei verlustfreien Strukturen und schlechter Energiebilanz in fp64 wiederholen \\
\key{split\_polarizations} & bool & \val{false} & 1D, $\psi=0$, isotrop: TE und TM getrennt (schneller) \\
\key{fft\_sampling} & bool & \val{false} & abgetastete Koeffizienten per FFT \\
\key{threads} & Objekt & auto & \key{scan}, \key{slices}, \key{blas} (Abschnitt~\ref{sec:threads}) \\
\key{reduction} & Text & \val{"sequential"} & Verkettung der Scheiben: \val{sequential} oder \val{tree} \\
\key{growth\_limit} & Zahl & 0 (auto) & Kriterium für die Stabilitätsscheiben \\
\key{checkpoint\_every} & Zahl & 1 & Gradienten: nur jedes $k$-te Zwischenprodukt speichern \\
\bottomrule
\end{tabularx}
\end{center}

\section{Ausgabe}
\begin{center}
\small
\begin{tabularx}{\linewidth}{lllY}
\toprule
Schlüssel & Typ & Standard & Bedeutung \\
\midrule
\key{output.dir} & Text & \val{"out"} & Ausgabeverzeichnis (relativ zum Arbeitsverzeichnis von \cmd{fgsolve}) \\
\key{output.efficiencies} & Text & \val{"efficiencies.csv"} & Name der Effizienzdatei \\
\key{output.materials} & Text & -- & optional: $n$ und $\varepsilon$ aller Materialien je Scan-Wellenlänge
(Abschnitt~\ref{sec:materials-csv}) \\
\bottomrule
\end{tabularx}
\end{center}

\section{Ein vollständiges Beispiel mit allen Grundblöcken}
\begin{lstlisting}[language=json]
{
  "dimension": "2D",
  "solver": "rcwa",
  "truncation": { "Nx": 6, "Nz": 6 },
  "period_um": { "x": 0.6, "z": 0.6 },
  "incidence": { "lambda_um": { "from": 0.55, "to": 0.75, "steps": 21 },
                 "phi_deg": 10, "psi_deg": 0, "delta_deg": [0, 90] },
  "materials": { "air": { "n": 1.0 }, "glass": { "n": 1.5 }, "si": { "n": [3.88, 0.02] } },
  "superstrate": "air",
  "substrate": "glass",
  "layers": [
    { "type": "polygon2d", "thickness_um": 0.15, "background": "air",
      "inclusions": [ { "material": "si",
                        "rectangle_um": { "center": [0.3, 0.3], "size": [0.3, 0.3] } } ] }
  ],
  "backend": "auto",
  "threads": { "scan": "auto", "slices": "auto", "blas": 1 },
  "output": { "dir": "out/si_pillars" }
}
\end{lstlisting}
Das sind 42 Rechnungen (21 Wellenlängen, zwei Polarisationen) an quadratischen Siliziumsäulen auf Glas.
